% Pregunta 659 (2008-julio-OEP2008, n.º 11): monopolio sin discriminar (IMg = CMg, q^M y p^M) frente a discriminación perfecta (demanda = CMg, q^C > q^M; la última unidad se vende a p = CMg < p^M). Autor: Víctor Gutiérrez Marcos
% A escala: p = 10 - q, IMg = 10 - 2q, CMg = 2.
\documentclass[tikz,border=6pt]{standalone}
\input{tikz-tcee}
\begin{document}
\begin{tikzpicture}[tcee,x=0.52cm,y=0.46cm]
  \ejes{$q$}{$p$}{11.5}{11}
  % Pérdida de eficiencia del monopolio que la discriminación perfecta recupera
  \fill[area2] (4,2) -- (4,6) -- (8,2) -- cycle;
  % Coste marginal constante
  \draw[curva3] (0,2) -- (10.6,2);
  \node[etiqueta,tceeverde,above] at (10.0,2.0) {$\mathit{CMg}$};
  % Demanda e ingreso marginal
  \draw[curva1] (0,10) -- (10,0);
  \node[etiqueta,tceemorado,right] at (9.3,0.9) {$D$};
  \draw[curva2] (0,10) -- (5,0);
  \node[etiqueta,tceeazul,left] at (4.2,1.2) {$\mathit{IMg}$};
  % Sin discriminar: IMg = CMg
  \draw[guia] (4,0) -- (4,6);
  \draw[guia] (0,6) -- (4,6);
  \node[left] at (0,6) {$p^M$};
  \node[below] at (4,-0.15) {$q^M$};
  \node[punto] at (4,2) {};
  \node[punto] at (4,6) {};
  \node[etiqueta,above right] at (4.05,6.05) {$M$};
  % Discriminación perfecta: D = CMg
  \draw[guia] (8,0) -- (8,2);
  \node[left] at (0,2) {$\mathit{CMg}$};
  \node[below] at (8,-0.15) {$q^C$};
  \node[punto] at (8,2) {};
  \node[etiqueta,above right] at (8.0,2.1) {$C$};
  % El precio marginal baja de p^M a CMg
  \draw[flecha,tceerojo] (-1.9,5.7) -- (-1.9,2.3);
  \node[etiqueta,tceerojo,align=center,font=\footnotesize,left] at (-2.0,4.0) {precio\\marginal};
  % Rótulo del área
  \node[etiqueta,tceerojo,align=left,font=\footnotesize] (r) at (7.6,6.6) {recuperada con\\discriminación\\perfecta};
  \draw[flecha,tceerojo,thin] (r.south) -- (5.4,3.3);
\end{tikzpicture}
\end{document}
